\chapter{Lezione 19 --- Riccati a regime, orizzonte infinito, MPC e programmazione dinamica in tempo continuo}

\section{Richiamo: controllo ottimo LQ a tempo discreto}

Consideriamo il sistema lineare discreto
\[
x_{k+1}=Ax_k+Bu_k,
\qquad x_0=\hat x,
\]
con funzionale di costo
\[
J=\sum_{k=0}^{N-1}\left(x_k^T V x_k+u_k^T P u_k\right)+x_N^T V_N x_N.
\]

Le matrici di costo soddisfano le ipotesi:
\[
V=V^T\succeq 0,
\qquad
P=P^T\succ 0,
\qquad
V_N=V_N^T\succeq 0.
\]

Dalle lezioni precedenti abbiamo dimostrato che il controllo ottimo ha la forma
\[
u_k^\circ=-H_k x_k,
\]
dove
\[
H_k=\left(P+B^T T_{k+1}B\right)^{-1}B^T T_{k+1}A,
\]
e il costo ottimo a partire dal passo \(k\) è
\[
J_k^\circ=x_k^T T_k x_k.
\]

La matrice \(T_k\) soddisfa la ricorsione all'indietro di Riccati:
\[
T_k
=
V
+
A^T
\left[
T_{k+1}
-
T_{k+1}B\left(P+B^T T_{k+1}B\right)^{-1}B^T T_{k+1}
\right]
A,
\]
con condizione terminale
\[
T_N=V_N.
\]

\section{Osservazioni sulla soluzione del problema LQ}

La soluzione del problema di controllo ottimo dipende dalla conoscenza del modello,
cioè delle matrici \(A\) e \(B\), e dalle matrici di peso \(V\), \(P\), \(V_N\).

Quindi, per poter implementare il controllo ottimo, è necessario:
\begin{itemize}
    \item disporre di un modello in equazioni di stato del sistema;
    \item scegliere le matrici di costo in base alle prestazioni desiderate;
    \item risolvere \emph{prima} dell'evoluzione del processo la ricorsione di Riccati.
\end{itemize}

In pratica, la ricorsione di Riccati viene risolta offline, partendo dall'istante finale
e andando all'indietro fino all'istante iniziale; una volta calcolati i guadagni
\(
H_k
\),
si può implementare online la legge di controllo
\[
u_k^\circ=-H_k x_k.
\]

\section{La ricorsione di Riccati e il concetto di regime}

Una particolarità della ricorsione di Riccati è che essa si risolve \emph{all'indietro}:
si parte da \(T_N\) e si ricavano \(T_{N-1},T_{N-2},\dots,T_0\).

Questo è diverso da ciò che accade per le normali equazioni di stato di un sistema stabile,
le cui soluzioni vengono tipicamente studiate in avanti nel tempo:
\[
x_0 \to x_1 \to x_2 \to \cdots
\]
e, nei sistemi asintoticamente stabili, si tende al regime per \(k\to+\infty\).

Nel caso di Riccati, essendo la ricorsione all'indietro, il concetto di regime va inteso
come \emph{punto fisso} della ricorsione:
\[
T_k=T_{k+1}=T.
\]

Sostituendo questa condizione nella ricorsione di Riccati si ottiene
l'\textbf{equazione algebrica di Riccati a tempo discreto}:
\[
T
=
V
+
A^T
\left[
T
-
TB\left(P+B^T T B\right)^{-1}B^T T
\right]A.
\]

Questa equazione non è più una ricorsione alle differenze, ma una
\textbf{equazione matriciale algebrica}.

\section{Orizzonte finito e orizzonte infinito}

L'orizzonte \(N\) ha un ruolo fondamentale.

\subsection{Orizzonte finito}

Se \(N\) è finito e fissato, si parla di \textbf{controllo ottimo su orizzonte finito}.
In questo caso:
\begin{itemize}
    \item bisogna calcolare una successione finita di matrici \(T_k\);
    \item si ottiene una successione finita di guadagni \(H_k\);
    \item il controllo ottimo è in generale \emph{tempo-variante}.
\end{itemize}

\subsection{Orizzonte infinito}

Se invece \(N\to+\infty\), si parla di \textbf{controllo ottimo su orizzonte infinito}.
In questo caso, sotto opportune ipotesi, la successione \(T_k\) converge a una matrice
costante \(T\), e quindi anche il guadagno converge a una matrice costante:
\[
H=\left(P+B^T T B\right)^{-1}B^T T A.
\]

Il controllo ottimo diventa allora
\[
u_k^\circ=-Hx_k,
\]
cioè una retroazione lineare a guadagno costante.

\subsection{Interpretazione}

Con orizzonte finito si ottimizza tenendo conto anche di come il sistema deve arrivare
al tempo finale fissato \(N\).

Con orizzonte infinito, invece, si rinuncia a un vincolo finale preciso:
si privilegia il comportamento complessivo su tempi lunghi.
In questo senso, l'ottimo a orizzonte infinito può essere meno raffinato sul transitorio
iniziale, ma produce una legge di controllo stazionaria molto più semplice da implementare.

\section{Implementazione pratica del controllo LQ}

Dal punto di vista implementativo:
\begin{enumerate}
    \item si costruisce il modello discreto del sistema;
    \item si scelgono le matrici \(V\), \(P\), \(V_N\);
    \item si risolve la ricorsione di Riccati all'indietro;
    \item si memorizzano i guadagni \(H_k\) oppure, nel caso a regime, il solo guadagno \(H\);
    \item durante il funzionamento si applica
    \[
    u_k^\circ=-H_kx_k
    \quad\text{oppure}\quad
    u_k^\circ=-Hx_k.
    \]
\end{enumerate}

\section{Dal controllo ottimo all'MPC}

L'idea dell'\textbf{MPC (Model Predictive Control)} nasce proprio dal controllo ottimo
su orizzonte finito.

In un regolatore MPC si fa, concettualmente, così:
\begin{enumerate}
    \item al tempo corrente si considera un orizzonte finito \(N\);
    \item si risolve il problema di controllo ottimo su tale orizzonte;
    \item si applica solo il primo tratto del controllo ottimo;
    \item si misura il nuovo stato;
    \item si sposta in avanti l'orizzonte e si ripete tutto da capo.
\end{enumerate}

Questa tecnica è detta anche \textbf{receding horizon control}.

\subsection{Idea operativa}

Negli appunti la logica è:
\begin{description}
    \item[Passo 1:] si linearizza il sistema attorno alle condizioni iniziali, si risolve
    Riccati su orizzonte \(N\), si applica il controllo LQ per un tempo \(H\ll N\);
    \item[Passo 2:] si osserva dove il sistema è arrivato, si rilinearizza attorno al nuovo
    punto di lavoro, si risolve di nuovo il problema LQ e si riapplica il controllo.
\end{description}

\section{Perché nasce l'MPC: impianti tempo-varianti}

Nella pratica gli impianti reali sono spesso:
\begin{itemize}
    \item non lineari;
    \item linearizzabili solo localmente;
    \item descrivibili da modelli che cambiano lungo la traiettoria.
\end{itemize}

Se il sistema evolve lungo una traiettoria, la linearizzazione locale può produrre matrici
tempo-varianti:
\[
A \to A_k,
\qquad
B \to B_k.
\]

In questo caso cambia anche la Riccati:
\[
T_k \to T_k \ \text{calcolata usando } A_k,B_k.
\]

Non esiste più, in generale, un'unica equazione algebrica di regime.
Bisogna invece risolvere una Riccati alle differenze con coefficienti variabili nel tempo.

Questo richiede di conoscere, a ogni istante, il modello locale del sistema.
Nella realtà questo è complesso, sia perché il modello è sempre approssimato,
sia perché possono agire disturbi e incertezze.

L'MPC è allora una soluzione pratica:
si ricalcola periodicamente il problema ottimo usando il modello locale più aggiornato.

\subsection{Condizione pratica}

Per fare funzionare questa strategia occorre che il processo non sia troppo veloce:
tra un aggiornamento e l'altro bisogna avere il tempo di
\begin{itemize}
    \item stimare o misurare lo stato;
    \item linearizzare il modello;
    \item risolvere il problema di Riccati;
    \item applicare il nuovo controllo.
\end{itemize}

\section{Programmazione dinamica nel caso a tempo continuo}

Passiamo ora al caso continuo, in forma più generale.

Consideriamo un impianto non lineare
\[
\dot x = F(x(t),u(t)),
\]
e supponiamo, per semplicità, che l'origine sia un equilibrio:
\[
F(0,0)=0.
\]

Vogliamo minimizzare, su orizzonte infinito, il funzionale di costo
\[
J=\int_{t_0}^{+\infty} \ell(x(\tau),u(\tau))\,d\tau.
\]

Siamo quindi interessati a trovare un controllo ottimo
\[
u^\circ(t)=\arg\min_{u(t)} J.
\]

\section{Richiamo del concetto di cost-to-go}

Nel tempo discreto avevamo introdotto il \textbf{cost-to-go}:
se un sistema si trova lungo una traiettoria ottima e arriva in un certo stato,
allora anche il tratto successivo della traiettoria deve restare ottimo.

Lo stesso concetto si usa qui.

Definiamo
\[
\psi(x_0)=J^\circ,
\]
cioè il costo ottimo associato alla condizione iniziale \(x_0\).

Più in generale, se al tempo \(t\) il sistema si trova nello stato \(x(t)\),
definiamo
\[
\psi(x(t))
=
\min_u \int_t^{+\infty}\ell(x(\tau),u(\tau))\,d\tau.
\]

\section{Scomposizione del costo su un intervallo breve}

Prendiamo un piccolo intervallo \(h>0\). Allora, per il principio di Bellman,
si può scrivere
\[
\psi(x(t))
=
\min_u
\left[
\int_t^{t+h}\ell(x(\tau),u(\tau))\,d\tau
+
\psi(x(t+h))
\right].
\]

Sottraendo \(\psi(x(t))\) da entrambi i membri:
\[
0
=
\min_u
\left[
\psi(x(t+h))-\psi(x(t))
+
\int_t^{t+h}\ell(x(\tau),u(\tau))\,d\tau
\right].
\]

Dividendo per \(h\):
\[
0
=
\min_u
\left[
\frac{\psi(x(t+h))-\psi(x(t))}{h}
+
\frac{1}{h}\int_t^{t+h}\ell(x(\tau),u(\tau))\,d\tau
\right].
\]

Passando al limite per \(h\to 0^+\), il secondo termine tende a
\[
\ell(x(t),u(t)),
\]
mentre il primo tende a
\[
\dot\psi(x(t)).
\]

Quindi otteniamo
\[
\min_u\left[\dot\psi(x(t))+\ell(x(t),u(t))\right]=0.
\]

\section{Regola della catena}

Poiché \(\psi\) dipende dallo stato e lo stato dipende dal tempo,
per la regola della catena:
\[
\dot\psi(x(t))
=
\nabla_x\psi(x(t))^T \dot x(t).
\]

Essendo
\[
\dot x(t)=F(x(t),u(t)),
\]
si ha
\[
\dot\psi(x(t))
=
\nabla_x\psi(x(t))^T F(x(t),u(t)).
\]

Sostituendo nella relazione precedente si ottiene la
\textbf{equazione di Bellman} (o, più precisamente, l'equazione di Hamilton--Jacobi--Bellman):
\[
\min_u
\left[
\nabla_x\psi(x)^T F(x,u)+\ell(x,u)
\right]
=0.
\]

Questa è una condizione necessaria che il costo ottimo deve soddisfare.

\section{Interpretazione dell'equazione di Bellman}

L'equazione
\[
\min_u
\left[
\nabla_x\psi(x)^T F(x,u)+\ell(x,u)
\right]
=0
\]
esprime il fatto che, in ogni stato \(x\), il controllo ottimo è quello che minimizza
la somma di:
\begin{itemize}
    \item costo istantaneo \(\ell(x,u)\);
    \item variazione istantanea del costo residuo \(\psi\) dovuta alla dinamica del sistema.
\end{itemize}

In altre parole, il controllo ottimo sceglie in ogni istante il miglior compromesso
tra costo immediato e costo futuro.

\section{Collegamento con il caso discreto}

Nel caso discreto avevamo:
\[
J_k^\circ(x_k)=\min_{u_k}\Bigl[h(x_k,u_k)+J_{k+1}^\circ(x_{k+1})\Bigr].
\]

Nel caso continuo, facendo tendere a zero il passo temporale,
questa relazione diventa la Bellman continua:
\[
\min_u
\left[
\nabla_x\psi(x)^T F(x,u)+\ell(x,u)
\right]=0.
\]

Quindi la programmazione dinamica continua è l'analogo infinitesimo
della ricorsione di Bellman discreta.

\section{Sintesi della lezione}

In questa lezione abbiamo visto che:

\begin{itemize}
    \item nel problema LQ discreto a orizzonte finito la soluzione si ottiene
    tramite la ricorsione all'indietro di Riccati;
    \item se la ricorsione raggiunge un punto fisso \(T\), si ottiene
    l'equazione algebrica di Riccati:
    \[
    T
    =
    V
    +
    A^T
    \left[
    T
    -
    TB(P+B^TTB)^{-1}B^TT
    \right]A;
    \]
    \item l'orizzonte finito produce guadagni tempo-varianti \(H_k\),
    mentre l'orizzonte infinito porta a un guadagno costante;
    \item l'MPC nasce dall'idea di risolvere ripetutamente problemi LQ su orizzonte finito,
    applicando ogni volta solo una parte del controllo ottimo;
    \item questa strategia è particolarmente utile per impianti non lineari o
    tempo-varianti, che vengono linearizzati localmente;
    \item nel caso continuo e non lineare, introducendo il cost-to-go
    \(\psi(x)\), si ottiene l'equazione di Bellman:
    \[
    \min_u\left[\nabla_x\psi(x)^T F(x,u)+\ell(x,u)\right]=0.
    \]
\end{itemize}

Questa equazione è il punto di partenza teorico della programmazione dinamica
nel continuo.